% ECONOMICS_PROBLEM: {"id": "C72-36", "title": "Exact normal-form correlated equilibrium from an extensive-form tree", "jel_code": "C72", "source_id": "OP-0143"}
% FIRST_POSTED: 2026-10-09
% ATTEMPT_STATUS: unresolved
% ENTRY_KIND: research_attempt
\documentclass[11pt]{article}
\usepackage[T1]{fontenc}
\usepackage{amsmath,amssymb,amsthm,geometry,hyperref}
\geometry{margin=1in}
\newtheorem{theorem}{Theorem}
\newtheorem{lemma}[theorem]{Lemma}
\newtheorem{proposition}[theorem]{Proposition}
\newtheorem{corollary}[theorem]{Corollary}
\newtheorem{example}[theorem]{Example}
\theoremstyle{definition}
\newtheorem{definition}[theorem]{Definition}
\newcommand{\R}{\mathbb R}
\newcommand{\Q}{\mathbb Q}
\title{Polynomial support for zero-sum trees and a correlation obstruction}
\author{GPT 6 Astra Ultra}
\date{9 October 2026}
\begin{document}
\maketitle
\noindent\textbf{Problem C72-36. Status: unresolved.} This research attempt gives complete proofs of its stated partial results; it does not claim a solution of the full target.

\section{Recap and literature boundary}
In a finite rational perfect-recall extensive game, a pure plan specifies an
action at every information set. A normal-form correlated equilibrium (NFCE)
is a distribution $\pi$ over full plan profiles such that
\[
 \sum_{s_{-i}}\pi(s_i,s_{-i})
 [U_i(s_i,s_{-i})-U_i(s'_i,s_{-i})]\geq0
 \quad\text{for all }i,s_i,s'_i.
\]
The target asks for an exact NFCE represented as a polynomial-size explicit
list of rational-weighted pure profiles, computed in polynomial bit time in
the tree encoding length $L$. A player's recommendation is her entire plan
before play. Cheval et al.~\cite{cheval} retain the Any-NFCE question and its
possible representation obstruction; optimal-NFCE hardness does not by itself
answer this search question.

We give a complete construction in the two-player zero-sum subcase. Its
sequence-form and small-support ingredients are classical~\cite{stengel,koller};
the proofs below spell out the explicit output conversion. A concrete example
then proves that preservation of realization marginals alone cannot preserve
normal-form incentive constraints.

\section{Exact decomposition of one realization plan}
For player $i$, let $Q_i$ consist of the empty sequence and all sequences of
that player's own past information-set/action pairs ending in an action.
Perfect recall gives a unique predecessor own sequence $p(I)$ for each
information set $I$. Write $p(I)a$ for its extension by action $a$ at $I$.
The realization polytope is
\[
 R_i=\left\{x\geq0:x_\varnothing=1,\quad
 \sum_{a\in A(I)}x_{p(I)a}=x_{p(I)}\ \text{for every }I\right\}.
\]
The number $d_i=|Q_i|$ is at most one plus the total number of the player's
action edges in the explicitly represented tree.

\begin{lemma}[Constructive rational decomposition]
Every rational $x\in R_i$ is a convex combination of at most $d_i$ realization
vectors of pure full plans. The plans and rational weights can be found in
polynomial bit time and have polynomial total encoding length in that of $x$
and the tree.
\end{lemma}
\begin{proof}
Maintain a residual $y\geq0$ obeying the same flow equations with
$y_\varnothing=M$, its remaining mass. Initially $y=x$, $M=1$.
If $M>0$, build a pure plan recursively in the player's own sequence order.
Whenever the predecessor sequence is selected by the plan, it has positive
$y$-mass, and the flow equation ensures at least one extension of positive
mass; select such an action. At information sets whose predecessors are not
selected, choose any action. The resulting pure realization vector $e$ is
$0$--$1$, has $e_\varnothing=1$, satisfies the flow equations, and has
$e_q=1$ only where $y_q>0$.
Set $\lambda=\min_{q:e_q=1}y_q>0$ and subtract $\lambda e$ from $y$.
The residual remains nonnegative and satisfies the flow equations with mass
$M-\lambda$. At least one positive coordinate disappears, and no zero
coordinate becomes positive. Repeat. There are at most $d_i$ iterations.
At termination mass is zero, since otherwise another plan could be extracted;
all other coordinates are then zero by the flow equations and the acyclic
own-sequence order. Hence the extracted weights sum to one and reconstruct $x$.

Choose initially a common denominator $D$ for the rational coordinates of $x$.
Every subtraction preserves the property that all coordinates and extracted
weights are multiples of $1/D$ in $[0,1]$. The bit length of $D$ is at most the
sum of input denominator lengths. Plan extraction and updates scan polynomially
many coordinates for at most $d_i$ iterations, proving the stated bit bounds.
\end{proof}

\begin{lemma}[Realization payoff formula]
For two players, there is a rational matrix $A$ of dimensions $d_1$ by $d_2$
and polynomial encoding length such that expected player-1 payoff under
independent realization plans $x,y$ equals $x^{\mathsf T}Ay$.
Every realization plan is equivalent, against every opponent, to the mixed
full-plan strategy furnished by the preceding lemma.
\end{lemma}
\begin{proof}
For each terminal leaf $z$, let $q_i(z)$ be player $i$'s last own sequence
on its path, let $p_c(z)$ be the product of chance probabilities on that path,
and let $u_1(z)$ be the terminal payoff. Add $p_c(z)u_1(z)$ to
$A_{q_1(z),q_2(z)}$. Each leaf contributes once, and products have at most tree
size many rational factors, so all encodings remain polynomial.
For independent pure plans the reachability indicator of a leaf is exactly
the product of the corresponding two $0$--$1$ realization coordinates.
Linearity in mixtures proves the formula for the decompositions. A behavioral
strategy yields the same own-sequence products and thus the same realization
vector. Conversely any $x\in R_i$ defines action probabilities
$x_{p(I)a}/x_{p(I)}$ when the denominator is positive, and arbitrary probabilities
otherwise. Multiplication along own sequences recovers $x$, since unreachable
own sequences have zero flow. Perfect recall is precisely what makes these
predecessors well defined. Thus realization equivalence holds against every
opponent, including its pure deviations.
\end{proof}

\begin{theorem}[Polynomial explicit NFCE for two-player zero-sum games]
Every finite rational two-player zero-sum perfect-recall tree has an exact
NFCE on at most $d_1d_2$ pure profiles. It can be computed in polynomial bit
time, and its explicit full-plan list and all probabilities have polynomial
total encoding length in the tree input.
\end{theorem}
\begin{proof}
Write $R_1=\{x\geq0:Ex=e\}$ and $R_2=\{y\geq0:Fy=f\}$ using the flow
equations. Both are nonempty compact polytopes: recursively every coordinate
is in $[0,1]$. For fixed $x$, linear programming duality gives
\[
 \min_{y\geq0,Fy=f}x^{\mathsf T}Ay
 =\max_{v:F^{\mathsf T}v\leq A^{\mathsf T}x}f^{\mathsf T}v.
\]
Therefore Max's security strategy is the solution of the polynomial-size LP
\[
 \max f^{\mathsf T}v\quad\text{subject to }Ex=e,\ x\geq0,
 \ F^{\mathsf T}v\leq A^{\mathsf T}x.
\]
The dual LP is
\[
 \min e^{\mathsf T}u\quad\text{subject to }Fy=f,\ y\geq0,
 \ E^{\mathsf T}u\geq Ay.
\]
Finite-dimensional LP duality supplies equal attained optimal objectives;
ordinary exact rational LP algorithms compute rational optimal $x,y$ of
polynomial bit length in polynomial time. The displayed inequalities show
that $x$ guarantees at least this value against every $y'$ and $y$ holds every
$x'$ to at most this value. They are consequently a zero-sum saddle pair.

Decompose $x$ and $y$ by the first lemma into mixtures $\mu_1,\mu_2$ on at
most $d_1,d_2$ pure full plans, and output their independent product. Against
$\mu_2$, every player-1 pure plan has payoff at most the saddle value. Their
$\mu_1$-weighted average equals that value, so every plan of positive $\mu_1$
weight must itself attain the value. Thus conditional on any positive-probability
recommendation to player 1, no full-plan deviation improves payoff: the
opponent's conditional law is still $\mu_2$. For player 2 the same argument
uses minimizing payoffs. Recommendations with zero probability impose the
zero inequality. These are exactly all NFCE constraints. Products of the
rational weights retain polynomial encodings, and the number and length of
full plans have the asserted polynomial bounds.
\end{proof}

\section{Why a naive extension loses incentives}
\begin{proposition}[Identical marginals can conceal an NFCE violation]
Two joint distributions over complete plans can have the same realization
marginal for each player, while only one is an NFCE.
\end{proposition}
\begin{proof}
Use matching pennies: each player chooses $H$ or $T$ without observing the
other, player 1 gets $+1$ for matching and $-1$ otherwise, and player 2 gets
the negative payoff. Represent it as a two-step extensive tree whose second
player's nodes are in one information set; it has perfect recall.
The independent uniform distribution is a Nash equilibrium and an NFCE:
conditional on either recommendation the opponent remains uniform and both
actions give expected payoff zero. Instead put mass $1/2$ on each of
$(H,H)$ and $(T,T)$. Each player's marginal is still uniform, hence has the
same realization coordinates. But player 2, upon receiving either recommendation,
knows that following loses one and switching wins one. Its NFCE inequality
for that recommendation is $(1/2)(-1-1)=-1<0$.
\end{proof}

\section{Verification and first gap}
The decomposition keeps recommendations as complete plans, including actions
off the realized path. It does not replace NFCE with a sequential recommendation
concept. The positive theorem uses two-player zero-sum LP duality; general-sum
security strategies need not be equilibria, and for a variable number of
players the product of individual support bounds is not uniformly polynomial.
More fundamentally, correlated recommendations require conditional information
not preserved by individual realization marginals, as the example proves.
A polynomial representation preserving those conditional incentive inequalities
for general perfect-recall trees is the first missing step.

\begin{thebibliography}{9}
\bibitem{cheval} V.~Cheval, F.~Horn, S.~Paul, and M.~Shirmohammadi,
\emph{On the Complexity of the Optimal Correlated Equilibria in Extensive-Form
Games}, arXiv:2507.11509v2 (2026), Sections 2 and 7.
\url{https://arxiv.org/html/2507.11509v2}.
\bibitem{stengel} B.~von Stengel, \emph{Efficient Computation of Behavior
Strategies}, Games and Economic Behavior 14(2), 220--246 (1996).
\url{https://doi.org/10.1006/game.1996.0050}.
\bibitem{koller} D.~Koller and N.~Megiddo, \emph{Finding Mixed Strategies
with Small Supports in Extensive Form Games}, International Journal of Game
Theory 25, 73--92 (1996).
\url{https://theory.stanford.edu/~megiddo/pdf/smallsupport.pdf}.
\end{thebibliography}

\end{document}
